\documentclass[10pt,a4paper]{amsart}
\usepackage[margin=19mm]{geometry}
\usepackage{amsmath,amssymb,mathtools}
\usepackage[scr=rsfs,cal=euler]{mathalfa}
\usepackage{xfrac}
\usepackage[unicode,hidelinks,bookmarks=false]{hyperref}
\newcommand{\ie}{i.e.\ }
\newcommand{\cf}{cf.\ }
\newcommand{\resp}{respectively\ }
\newcommand{\Subsection}{\subsection}
\providecommand{\nobreakdash}{\nobreak\hbox{-}}
\setlength{\emergencystretch}{2em}
\multlinegap0pt
\usepackage{enumerate}
\usepackage{amssymb}
\newtheorem{lemma}{Lemma}
\title{Order extreme points, solid convex hull, and lattice analogues of Krein-Milman Theorem}
\author{Anastasiia Ianina, Timur Oikhberg, Mary Angelica Tursi}
\date{Source: arXiv:2606.19415v1 (CC BY 4.0)}
\begin{document}
\maketitle

\noindent\emph{Source and licence.} Order extreme points, solid convex hull, and lattice analogues of Krein-Milman Theorem. Version arXiv:2606.19415v1 (CC BY 4.0), distributed by arXiv under the Creative Commons Attribution 4.0 International licence, http://creativecommons.org/licenses/by/4.0/, as recorded in the arXiv licence metadata for this exact version and shown on the abstract page https://arxiv.org/abs/2606.19415v1. Full licence notice: \texttt{licenses/arxiv-2606.19415-CC-BY-4.0.txt}.\par\smallskip

\noindent\textit{Editorial scope.} Counted unit: the article's lemma l:shadow\_compact (source0.tex lines 475--477). The article's complete original proof of it is delivered with it (source0.tex lines 479--483, one \texttt{proof} environment of 5 lines). No mathematics is written, repaired or re-derived: the statement and the proof are the article's own lines, in the article's own notation, and no retained line is substituted or re-flowed.  No further source-local context is needed: every cross-reference of every retained line resolves inside the two delivered spans.  Nothing was omitted from any retained span, so the package carries no omission record.  No retained line contains a citation, so the wrapper carries no bibliography and no bibliography entry.  No retained line contains a figure environment, an image-inclusion command or drawing code, so no image is delivered and the package declares no asset dependency.  Editorial formatting only, in addition: an AMS \texttt{amsart} preamble that restates the article's own declarations and macros used by the retained lines, namely the environments \texttt{lemma} on separate counters, together with the standard packages those lines need.  The retained environments print in this document as Lemma 1 in order of appearance, whereas the article prints the same items with the numbering of its own full document.  The complete native source of the article as distributed in the arXiv e-print for this exact version is bundled unchanged as \texttt{sources/203141/source0.tex}.
\begin{lemma}\label{l:shadow_compact}
If $A \subset X_+$ is compact, then $A_s$ is closed.
\end{lemma}

\begin{proof}
Suppose the net $(x_\alpha)_{\alpha \in {\mathcal{I}}} \subset A_s$ converges to $x$, and show that $x \in A_s$ as well.
By the definition of $A_s$, for any $\alpha$ there exists $a_\alpha \in A$ so that $x_\alpha \leq a_\alpha$.
By compactness, the net ${\mathcal{I}}$ has a cofinite subnet ${\mathcal{J}}$ so that $(a_\beta)_{\beta \in {\mathcal{J}}}$ converges to some $a \in A$. Then $a \geq x$, hence $x \in A_s$. 
\end{proof}

\end{document}
